\subsection{扩张模的定义与基本性质}

对于每个 A-模 E，我们用 \(p_{\mathbf{E}} : \mathbf{L}(\mathbf{E}) \to \mathbf{E}\) （相应地， \(e_{\mathbf{E}} : \mathbf{E} \to \mathbf{I}(\mathbf{E})\) ）表示典范的自由（相应地，内射）分解，参见 X，第 50 页（相应地，第 52 页）。

定义1. --- 设 M 和 N 是两个 A-模。称 N 由 M 扩张的扩张模为分次 k-模

\[
\operatorname{Ext} _ {\mathbf {A}} (\mathbf {M}, \mathbf {N}) = \mathrm{H} \left(\operatorname{Homgr} _ {\mathbf {A}} (\mathrm{L} (\mathbf {M}), \mathrm{I} (\mathbf {N}))\right).\tag{4}
\]

Ext \(_{A}\) (M, N) 的齐次分量记为

\[
\operatorname{Ext} _ {\mathbf {A}} ^ {n} (\mathbf {M}, \mathbf {N}) = \mathrm{H} ^ {n} \left(\operatorname{Homgr} _ {\mathbf {A}} (\mathbf {L} (\mathbf {M}), \mathrm{I} (\mathbf {N}))\right).\tag{5}
\]

由于 L(M)（相应地，I(N)）右零（相应地，左零），我们有

\[
\operatorname{Ext} _ {\mathbf {A}} ^ {n} (\mathbf {M}, \mathbf {N}) = 0 \quad \text { for } \quad n <   0.\tag{6}
\]

注记1. --- 我们将在下文（X，第 107 页，命题 6）看到 Ext \(_{A}\) (M, N) 这些模的有限性性质。例如，若 A 是交换诺特环，且 M 与 N 是有限生成的 A-模，则每个 A-模 Ext \(_{A}^{n}\) (M, N) 是有限生成的。

设 \(f: \mathbf{M}' \to \mathbf{M}\) 和 \(g: \mathbf{N} \to \mathbf{N}'\) 为 A-模的同态，我们设

\[
\operatorname{Ext} _ {\mathbf {A}} (f, g) = \mathrm{H} \left(\operatorname{Homgr} _ {\mathbf {A}} (\mathrm{L} (f), \mathrm{I} (g))\right);
\]

它是分次 k-模的 0 次同态

\[
\operatorname{Ext} _ {\mathbf {A}} (f, g): \operatorname{Ext} _ {\mathbf {A}} (\mathbf {M}, \mathbf {N}) \rightarrow \operatorname{Ext} _ {\mathbf {A}} \left(\mathbf {M} ^ {\prime}, \mathbf {N} ^ {\prime}\right),
\]

其齐次分量记为

\[
\operatorname{Ext} _ {\mathbf {A}} ^ {n} (f, g): \operatorname{Ext} _ {\mathbf {A}} ^ {n} (\mathbf {M}, \mathbf {N}) \rightarrow \operatorname{Ext} _ {\mathbf {A}} ^ {n} \left(\mathbf {M} ^ {\prime}, \mathbf {N} ^ {\prime}\right).
\]

由 X 第 82 页命题 1，典范同态

\[
\lambda^ {0} (\mathrm{L} (\mathrm{M}), \mathrm{I} (\mathrm{N})): \mathrm{H} ^ {0} \left(\operatorname{Homgr} _ {\mathrm{A}} (\mathrm{L} (\mathrm{M}), \mathrm{I} (\mathrm{N}))\right)\rightarrow \operatorname{Hom} _ {\mathrm{A}} \left(\mathrm{H} _ {0} (\mathrm{L} (\mathrm{M})), \mathrm{H} ^ {0} (\mathrm{I} (\mathrm{N}))\right)
\]

是双射的；利用同构 \(\mathbf{M} \to \mathrm{H}_{0}(\mathrm{L}(\mathbf{M}))\) 和 \(\mathrm{H}^{0}(\mathrm{I}(\mathbf{N})) \to \mathbf{N}\) ，我们推出一个同构，称之为典范的

\[
\lambda_ {\mathbf {M}, \mathbf {N}}: \operatorname{Ext} _ {\mathbf {A}} ^ {0} (\mathbf {M}, \mathbf {N}) \rightarrow \operatorname{Hom} _ {\mathbf {A}} (\mathbf {M}, \mathbf {N}).\tag{7}
\]

我们总是通过这个同构把 \(\operatorname{Ext}_{\mathbf{A}}^{0}(\mathbf{M},\mathbf{N})\) 与 \(\operatorname{Hom}_{\mathbf{A}}(\mathbf{M},\mathbf{N})\) 等同。于是 \(k\) -线性映射 \(\operatorname{Ext}_{\mathbf{A}}^{0}(f,g)\) 等同于 \(\operatorname{Hom}_{\mathbf{A}}(f,g)\) .

注记2. --- 复形态射

\[
\operatorname{Homgr} _ {\mathbf {A}} \left(p _ {\mathbf {M}}, e _ {\mathbf {N}}\right): \operatorname{Hom} _ {\mathbf {A}} (\mathbf {M}, \mathbf {N}) \rightarrow \operatorname{Homgr} _ {\mathbf {A}} \left(\mathbf {L} (\mathbf {M}), \mathrm{I} (\mathbf {N})\right)
\]

在0次同调上诱导同构

\[
\lambda_ {\mathbf {M}, \mathbf {N}} ^ {- 1}: \operatorname{Hom} _ {\mathbf {A}} (\mathbf {M}, \mathbf {N}) \rightarrow \operatorname{Ext} _ {\mathbf {A}} ^ {0} (\mathbf {M}, \mathbf {N})
\]

\(\lambda_{M,N}\) 的逆.

我们有 \(L(1_{\mathbf{M}}) = 1_{L(\mathbf{M})}\) , \(I(1_{\mathbf{N}}) = 1_{I(\mathbf{N})}\) ，因此通过过渡到同调

\[
\operatorname{Ext} _ {\mathbf {A}} \left(1 _ {\mathbf {M}}, 1 _ {\mathbf {N}}\right) = 1 _ {\operatorname{Ext} _ {\mathbf {A}} (\mathbf {M}, \mathbf {N})}.\tag{8}
\]

若 \(f': M'' \to M'\) 和 \(g': N' \to N''\) 是 A-模的同态，我们有 \(L(f \circ f') = L(f) \circ L(f')\) 和 \(I(g' \circ g) = I(g') \circ I(g)\) ，因此

\[
\operatorname{Ext} _ {\mathbf {A}} \left(f \circ f ^ {\prime}, g ^ {\prime} \circ g\right) = \operatorname{Ext} _ {\mathbf {A}} \left(f ^ {\prime}, g ^ {\prime}\right) \circ \operatorname{Ext} _ {\mathbf {A}} (f, g).\tag{9}
\]

考虑 k-复形的态射

\[
\operatorname{Homgr} _ {\mathbf {A}} (\mathrm{L} (\mathrm{M}), \mathrm{N}) \xrightarrow {\operatorname{Homgr} _ {\mathbf {A}} (1 , e _ {\mathrm{N}})} \operatorname{Homgr} _ {\mathbf {A}} (\mathrm{L} (\mathrm{M}), \mathrm{I} (\mathrm{N})) \xleftarrow {\operatorname{Homgr} _ {\mathbf {A}} (p _ {\mathrm{M}}, 1)} \operatorname{Homgr} _ {\mathbf {A}} (\mathrm{M}, \mathrm{I} (\mathrm{N})),
\]

以及它们在同调中诱导的同态

\[
\mathrm{H} \left(\operatorname{Homgr} _ {\mathrm{A}} (\mathrm{L} (\mathrm{M}), \mathrm{N})\right) \xrightarrow {\varphi_ {\mathrm{M}} (\mathrm{N})} \operatorname{Ext} _ {\mathrm{A}} (\mathrm{M}, \mathrm{N}) \xleftarrow {\overline {{\varphi}} _ {\mathrm{N}} (\mathrm{M})} \mathrm{H} \left(\operatorname{Homgr} _ {\mathrm{A}} (\mathrm{M}, \mathrm{I} (\mathrm{N}))\right).
\]

根据 X 第 86 页命题 4， \(\operatorname{Homgr}_{\mathbf{A}}(1, e_{\mathbf{N}})\) 和 \(\operatorname{Homgr}_{\mathbf{A}}(p_{\mathbf{M}}, 1)\) 是拟同构，因此：

命题5. --- k-同态

\[
\begin{array}{l l} & \varphi_ {\mathbf {M}} (\mathbf {N}): \mathrm{H} (\mathrm{Homgr} _ {\mathbf {A}}   (\mathrm{L} (\mathbf {M}), \mathrm{N})) \to \mathrm{Ext} _ {\mathbf {A}}   (\mathbf {M}, \mathbf {N}) \\ 	ext{and} & \overline {{\varphi}} _ {\mathbf {N}} (\mathbf {M}): \mathrm{H} (\mathrm{Homgr} _ {\mathbf {A}}   (\mathbf {M}, \mathrm{I} (\mathbf {N}))) \to \mathrm{Ext} _ {\mathbf {A}}   (\mathbf {M}, \mathbf {N}) 	ext{are bijective.} \end{array}
\]

推论. --- 若 M 是投射的（相应地，若 N 是内射的），则 Ext \(_{A}^{i}\) (M, N) = 0 （对于 i \textgreater{} 0）.

事实上，

\[
\operatorname{Homgr} _ {\mathbf {A}} \left(1, e _ {\mathbf {N}}\right): \operatorname{Hom} _ {\mathbf {A}} (\mathbf {M}, \mathbf {N}) \rightarrow \operatorname{Homgr} _ {\mathbf {A}} (\mathbf {M}, \mathrm{I} (\mathbf {N}))
\]

（相应地，Homgr \(_{A}\) ( \(p_{M}\) ，1)：Hom \(_{A}\) (M, N) → Homgr \(_{A}\) (L(M), N)) 于是是一个拟同构（X，第 86 页，命题4），由此得出结论。

注记。--- 3) 若 \(g: \mathbf{N} \to \mathbf{N}'\) 是 A-模的同态，则

\[
\operatorname{Homgr} _ {\mathbf {A}} \left(1 _ {\mathrm{L} (\mathbf {M})}, g\right) \circ \operatorname{Homgr} _ {\mathbf {A}} \left(1 _ {\mathrm{L} (\mathbf {M})}, e _ {\mathbf {N}}\right) = \operatorname{Homgr} _ {\mathbf {A}} \left(1 _ {\mathrm{L} (\mathbf {M})}, e _ {\mathbf {N} ^ {\prime}}\right) \circ \operatorname{Homgr} _ {\mathbf {A}} \left(1 _ {\mathrm{L} (\mathbf {M})}, \mathrm{I} (g)\right)
\]

因此图表

\[
\begin{array}{c} \mathrm{H} (\operatorname{Homgr} _ {\mathbf {A}} (\mathrm{L} (\mathbf {M}), \mathrm{N})) \xrightarrow {\varphi_ {\mathbf {M}} (\mathrm{N})} \operatorname{Ext} _ {\mathbf {A}} (\mathbf {M}, \mathbf {N}) \\ \mathrm{H} (\operatorname{Homgr} _ {\mathbf {A}} (1 _ {\mathrm{L} (\mathbf {M})}, g)) \Bigg | _ {\downarrow} \qquad \qquad \qquad \qquad \qquad \qquad \qquad \qquad \qquad \qquad \qquad \qquad \qquad \qquad \qquad \qquad \qquad \qquad \qquad \qquad \qquad \qquad \qquad \qquad \qquad \qquad \qquad \qquad \qquad \qquad \qquad \qquad \qquad \qquad \\ \mathrm{H} (\operatorname{Homgr} _ {\mathbf {A}} (\mathrm{L} (\mathbf {M}), \mathrm{N} ^ {\prime})) \xrightarrow {\varphi_ {\mathbf {M}} (\mathrm{N} ^ {\prime})} \operatorname{Ext} _ {\mathbf {A}} (\mathbf {M}, \mathrm{N} ^ {\prime}) \end{array}
\]

是交换的。

\begin{enumerate}
\def\labelenumi{\arabic{enumi})}
\setcounter{enumi}{3}
\tightlist
\item
  同样地，若 \(f: \mathbf{M}' \to \mathbf{M}\) 是 A-模的同态，则图表
\end{enumerate}

\[
\begin{array}{c c c} \mathrm{H} (\mathrm{Homgr} _ {\mathrm{A}} (\mathrm{M}, \mathrm{I} (\mathrm{N}))) & \xrightarrow {\overline {{\varphi}} _ {\mathrm{N}} (\mathrm{M})} & \mathrm{Ext} _ {\mathrm{A}} (\mathrm{M}, \mathrm{N}) \\ \mathrm{H} (\mathrm{Homgr} _ {\mathrm{A}} (f, 1 _ {\mathrm{I} (\mathrm{N})})) \Big \downarrow & & \Big \downarrow \mathrm{Ext} _ {\mathrm{A}} (f, 1 _ {\mathrm{N}}) \\ \mathrm{H} (\mathrm{Homgr} _ {\mathrm{A}} (\mathrm{M} ^ {\prime}, \mathrm{I} (\mathrm{N}))) & \xrightarrow {\overline {{\varphi}} _ {\mathrm{N}} (\mathrm{M} ^ {\prime})} & \mathrm{Ext} _ {\mathrm{A}} (\mathrm{M} ^ {\prime}, \mathrm{N}) \end{array}
\]

是交换的。

命题 6. --- 映射 (f, g) ↦ Ext\_A (f, g) :

\[
\operatorname{Hom} _ {\mathbf {A}} \left(\mathrm{M} ^ {\prime}, \mathrm{M}\right) \times \operatorname{Hom} _ {\mathbf {A}} \left(\mathrm{N}, \mathrm{N} ^ {\prime}\right)\rightarrow \operatorname{Homgr} _ {k} \left(\operatorname{Ext} _ {\mathbf {A}} (\mathrm{M}, \mathrm{N}), \operatorname{Ext} _ {\mathbf {A}} \left(\mathrm{M} ^ {\prime}, \mathrm{N} ^ {\prime}\right)\right)
\]

是 k-双线性的。

设 \(f \in \operatorname{Hom}_{\mathbf{A}}(\mathbf{M}', \mathbf{M})\) , \(g_1, g_2 \in \operatorname{Hom}_{\mathbf{A}}(\mathbf{N}, \mathbf{N}')\) , \(\lambda_1, \lambda_2 \in k\) ；则态射 \(\operatorname{Homgr}_{\mathbf{A}}(\mathbf{L}(f), \lambda_1 g_1 + \lambda_2 g_2)\) 和 \(\lambda_1 \operatorname{Homgr}_{\mathbf{A}}(\mathbf{L}(f), g_1) + \lambda_2 \operatorname{Homgr}_{\mathbf{A}}(\mathbf{L}(f), g_2)\) （从 \(\operatorname{Homgr}_{\mathbf{A}}(\mathbf{L}(\mathbf{M}), \mathbf{N})\) 到 \(\operatorname{Homgr}_{\mathbf{A}}(\mathbf{L}(\mathbf{M}), \mathbf{N}')\) ）重合；因此，由命题 5 和注记 3，

\[
\operatorname{Ext} _ {\mathbf {A}} \left(f, \lambda_ {1} g _ {1} + \lambda_ {2} g _ {2}\right) = \lambda_ {1} \operatorname{Ext} _ {\mathbf {A}} \left(f, g _ {1}\right) + \lambda_ {2} \operatorname{Ext} _ {\mathbf {A}} \left(f, g _ {2}\right).\tag{10}
\]

我们对映射 \(f\mapsto\mathrm{Ext}_{\mathbf{A}}(f,g)\) 同样地进行推理.

推论. --- 设 \(\lambda \in k\) 。如果 \(\lambda\) 零化 M 或 N，则它零化 Ext \(_{A}\) (M, N)。

事实上， \(\lambda 1_{\mathrm{Ext}(\mathbf{M},\mathbf{N})} = \mathrm{Ext}(\lambda 1_{\mathbf{M}},1_{\mathbf{N}}) = \mathrm{Ext}(1_{\mathbf{M}},\lambda 1_{\mathbf{N}})\) .

命题7. --- 设 I 和 J 为两个集合， \((\mathbf{M}_{\alpha})_{\alpha \in I}\) 和 \((\mathbf{N}_{\beta})_{\beta \in J}\) 为两个 A-模族；同态 \(\operatorname{Ext}_{\mathbf{A}}\left(\bigoplus_{\alpha \in I} \mathbf{M}_{\alpha}, \prod_{\beta \in J} \mathbf{N}_{\beta}\right) \to \prod_{\beta \in J, \alpha \in I} \operatorname{Ext}_{\mathbf{A}}(\mathbf{M}_{\alpha}, \mathbf{N}_{\beta})\) （它由典范同态 \(\mathbf{M}_{\alpha} \to \oplus \mathbf{M}_{\alpha}\) 和 \(\Pi \mathbf{N}_{\beta} \to \mathbf{N}_{\beta}\) 推出）是双射。

只需证明对每个 A-模 M（相应地，N），同态

\[
\operatorname{Ext} \left(\mathrm{M}, \Pi \mathrm{N} _ {\beta}\right)\rightarrow \Pi \operatorname{Ext} \left(\mathrm{M}, \mathrm{N} _ {\beta}\right) \quad (\text { resp.   Ext } (\oplus \mathrm{M} _ {\alpha}, \mathrm{N}) \rightarrow \Pi \operatorname{Ext} \left(\mathrm{M} _ {\alpha}, \mathrm{N}\right))
\]

是双射的。现在这由前述内容、X 的命题1（第 28 页）以及典范同构 \(\operatorname{Homgr}_{\mathbf{A}}(\mathbf{L}(\mathbf{M}), \Pi \mathbf{N}_{\beta}) \to \Pi \operatorname{Homgr}_{\mathbf{A}}(\mathbf{L}(\mathbf{M}), \mathbf{N}_{\beta})\) 和 \(\operatorname{Homgr}_{\mathbf{A}}(\oplus \mathbf{M}_{\alpha}, \mathrm{I}(\mathbf{N})) \to \Pi \operatorname{Homgr}_{\mathbf{A}}(\mathbf{M}_{\alpha}, \mathrm{I}(\mathbf{N}))\) 得出.

注记5.——设 P 与 Q 为两个右 A-模。定义 Ext \(_{A}\) (P, Q) 为

\[
\begin{array}{r l} \operatorname{Ext} _ {\mathbf {A}} (\mathrm{P}, \mathrm{Q}) & = \mathrm{H} (\operatorname{Homgr} _ {\mathbf {A}} (\mathrm{L} (\mathrm{P}), \mathrm{I} (\mathrm{Q}))) = \mathrm{H} (\operatorname{Homgr} _ {\mathbf {A} ^ {\circ}} (\mathrm{L} (\mathrm{P} ^ {\circ}), \mathrm{I} (\mathrm{Q} ^ {\circ}))) = \\ & = \operatorname{Ext} _ {\mathbf {A} ^ {\circ}} (\mathrm{P} ^ {\circ}, \mathrm{Q} ^ {\circ}). \end{array}
\]

因此，本段中的所有定义与命题均适用于右 A-模，只需把它们视为环 A° 上的左模。

